\documentclass[10pt,a4paper]{amsart}
\usepackage[margin=19mm]{geometry}
\usepackage{amsmath,amssymb,mathtools}
\usepackage[scr=rsfs,cal=euler]{mathalfa}
\usepackage{xfrac}
\usepackage[unicode,hidelinks,bookmarks=false]{hyperref}
\newcommand{\ie}{i.e.\ }
\newcommand{\cf}{cf.\ }
\newcommand{\resp}{respectively\ }
\newcommand{\Subsection}{\subsection}
\providecommand{\nobreakdash}{\nobreak\hbox{-}}
\setlength{\emergencystretch}{2em}
\multlinegap0pt
\usepackage{amssymb}
\newtheorem{lem}{Lemma}
\newtheorem{thm}{Theorem}
\newtheorem{prop}{Proposition}
\newtheorem{pf}{Pf}
\title{Convergence Analysis of Natural Power Method and Its Applications to Control}
\author{Daiki Tsuzuki, Kentaro Ohki}
\date{Source: arXiv:2512.21469v2 (CC BY 4.0)}
\begin{document}
\maketitle

\noindent\emph{Source and licence.} Convergence Analysis of Natural Power Method and Its Applications to Control. Version arXiv:2512.21469v2 (CC BY 4.0), distributed by arXiv under the Creative Commons Attribution 4.0 International licence, http://creativecommons.org/licenses/by/4.0/, as recorded in the arXiv licence metadata for this exact version and shown on the abstract page https://arxiv.org/abs/2512.21469v2. Full licence notice: \texttt{licenses/arxiv-2512.21469-CC-BY-4.0.txt}.\par\smallskip

\noindent\textit{Editorial scope.} Counted unit: the article's prop prop:ensure\_stationary\_point (source0.tex lines 359--369). The article's complete original proof of it is delivered with it (source0.tex lines 371--393, one \texttt{proof} environment of 23 lines). No mathematics is written, repaired or re-derived: the statement and the proof are the article's own lines, in the article's own notation, and no retained line is substituted or re-flowed.  Retained uncounted as the source-local context the proof needs: lines 59--62; lines 139--150; lines 218--227; lines 338--343.  Nothing was omitted from any retained span, so the package carries no omission record.  No retained line contains a citation, so the wrapper carries no bibliography and no bibliography entry.  No retained line contains a figure environment, an image-inclusion command or drawing code, so no image is delivered and the package declares no asset dependency.  Editorial formatting only, in addition: an AMS \texttt{amsart} preamble that restates the article's own declarations and macros used by the retained lines, namely the environments \texttt{lem}, \texttt{thm}, \texttt{prop}, \texttt{pf} on separate counters, together with the standard packages those lines need.  The retained environments print in this document as Lemma 1, Theorem 1, Proposition 1, Pf 1 in order of appearance, whereas the article prints the same items with the numbering of its own full document.  The complete native source of the article as distributed in the arXiv e-print for this exact version is bundled unchanged as \texttt{sources/191662/source0.tex}.
\begin{align}
    U[k+1] =& A U[k] (U[k]^{\top} A^{\top} A U[k]) ^{-1/2}, 
    \label{eq:natural_power_method}
\end{align}

\begin{lem}\label{lemma:discrete_norank_deficit}
Let $m\geq r$. If the initial condition $U[0]$ is in 
\begin{align}
    \mathcal{V}_{m,r}^{\rm (d)} := \Bigg{\{} \Psi [A] & \begin{bmatrix} K_{m,r} \\ K_{\perp} \end{bmatrix}  \in \mathrm{St}(r,n) \ \Bigg{|} \ 
    K_{m,r} \in \mathbb{C}^{m\times r},
     \nonumber \\ &
    K_{\perp} \in \mathbb{C}^{(n-m) \times r}, 
    \ \mathrm{rank}(K_{m,r}) =r \Bigg{\}}
    \label{eq:discrete_Oja_domain}
\end{align}
and the eigenvalues satisfy $| \lambda _{m}[A]| > | \lambda _{m+1}[A]|$, then $U[k] \in \mathcal{V}_{m,r}^{\rm (d)}$ and $U[k]^{\top}A^{\top}AU[k]$ is non-singular $\forall k\geq 0$.
\end{lem}

\begin{thm}\label{thm:discrete_eigenvector_subspace}
Suppose that for a given $A\in \mathbb{R}^{n\times n}$, there exists $m \in \{r, \dots ,n-1\}$ such that $| \lambda _{m} [A] | > | \lambda _{m+1} [A] | $. 
If the initial condition $U[0] \in \mathcal{V}_{m,r}^{\rm (d)}$, where $\mathcal{V}_{m,r}^{\rm (d)}$ is defined in Eq. \eqref{eq:discrete_Oja_domain}, then the solution $U[k]$ of \eqref{eq:natural_power_method} converges to the invariant set
\begin{align}
    \mathcal{U}_{m,r}^{\rm (d)} := \Bigg{\{} \Psi [A] \begin{bmatrix} K_{m,r} \\ 0_{n-m,r} \end{bmatrix}  \in & \mathrm{St}(r,n)  \ \Bigg{|}  \ K_{m,r} \in \mathbb{C}^{m\times r},
    \nonumber \\ &
    \ \mathrm{rank}(K_{m,r}) =r \Bigg{\}}.
    \label{eq:discrete_equillibriums}
\end{align}
\end{thm}

\begin{align}
    A \bar{U}_{r} = & \Psi _{r} \Lambda _{r} K_{r}
    = 
    \Psi _{r} K_{r} K_{r}^{-1} \Lambda _{r} K_{r} =  \bar{U} _{r}A_{\bar{U}_{r}}, 
    \label{eq:exponential_property}
\end{align}

\begin{prop}\label{prop:ensure_stationary_point}
Assume that $A$ satisfies $| \lambda _{r}[A] | > | \lambda _{r+1}[A] |$. Then, the sequence generated by the following algorithm converges to an element of $\mathcal{U}_{r}^{\rm (d)}$ if $\hat{U}[0] \in \mathcal{V}_{r} ^{\rm (d)}$ and $\hat{U}[k]^{\top}A \hat{U}[k]$ is non-singular for all $k\geq 0$:
\begin{align}
    \hat{U}[k+1] =& A \hat{U}[k] (\hat{U}[k]^{\top} A^{\top} A \hat{U}[k]) ^{-1/2} 
    \nonumber \\
    & \times (\hat{U}[k]^{\top} A^{\top} \hat{U}[k] \hat{U}[k]^{\top} A \hat{U}[k]) ^{-1/2} 
    \nonumber \\ & \times
    \hat{U}[k]^{\top} A^{\top} \hat{U}[k].
    \label{eq:ensure_stationary_point}
\end{align}
\end{prop}

\begin{pf}
Note that the difference between \eqref{eq:natural_power_method} and \eqref{eq:ensure_stationary_point} is that \eqref{eq:ensure_stationary_point} includes an additional multiplication $W[k] := (\hat{U}[k]^{\top} A^{\top} \hat{U}[k] \hat{U}[k]^{\top} A \hat{U}[k]) ^{-1/2} (\hat{U}[k]^{\top} A^{\top} \hat{U}[k]) \in \mathbb{R}^{r\times r}$, which is an orthogonal matrix.

Under the assumptions and the arguments in the proof of Lemma \ref{lemma:discrete_norank_deficit}, the solution $\hat{U}[k]$ of the algorithm \eqref{eq:ensure_stationary_point} always satisfies the non-singularity of $\hat{U}[k]^{\top} A^{\top} A \hat{U}[k]$.  
Additionally, since $W[k]$ is an orthogonal matrix, the same argument as in Theorem \ref{thm:discrete_eigenvector_subspace} implies that $U[k]$ of \eqref{eq:ensure_stationary_point} converges to $\mathcal{U}_{r}^{\rm (d)}$. 
Thus, the remainder of the proof focuses on stationarity. 

Let $\bar{U}_{r} \in \mathcal{U}_{r}^{\rm (d)}$. Then there exists a non-singular matrix $K_{r} \in \mathbb{C}^{r\times r}$ such that $\bar{U}_{r} = \Psi _{r} K_{r}$, where $\Psi _{r} := [\psi _{1}[A] ,\dots ,\psi _{r}[A]]$.  
From Eq. \eqref{eq:exponential_property}, $A \bar{U} _{r} = \bar{U}_{r} A_{\bar{U}_{r}}$, where $A_{\bar{U}_{r}} := \bar{U}_{r}^{\top}A\bar{U}_{r} \in \mathbb{R}^{r\times r}$, resulting in $\bar{U}_{r}^{\top} A^{\top} A \bar{U}_{r} = A_{\bar{U}_{r}}^{\top} A_{\bar{U}_{r}}$. 
On the other hand,
\begin{align*}
    &(\bar{U}_{r}^{\top} A^{\top} A \bar{U}_{r} ) ^{-1/2} 
    (\bar{U}_{r}^{\top} A^{\top} \bar{U}_{r} \bar{U}_{r}^{\top} A \bar{U}_{r} ) ^{-1/2} 
    \bar{U}_{r}^{\top} A \bar{U}_{r}
    \\
    =&
    (A_{\bar{U}_{r}}^{\top} A_{\bar{U}_{r}}) ^{-1} A_{\bar{U}_{r}}^{\top} = A_{\bar{U}_{r}}^{-1}
\end{align*}
holds. 
Therefore, if $\hat{U}[0] = \bar{U}_{r}  \in \mathcal{U}_{r}^{\rm (d)}$, then
    $\hat{U}[1] = A\bar{U}_{r} A_{\bar{U}_{r}}^{-1} = \bar{U}_{r} = \hat{U}[0]$.
This result concludes that the sequence generated by \eqref{eq:ensure_stationary_point} converges to an element of $\mathcal{U}_{r}^{\rm (d)}$.  
\end{pf}

\end{document}
